\FloatBarrier
\newpage
\section{导数的应用}
\label{sec:calculus-applications}

中值定理把导数与函数增量联系起来，因而可以由导数研究区间上的单调性，再确定极值、最值和方程的根。含参数的问题也遵循这个次序，只是需要继续判断这些结构在什么参数处发生改变。

\subsection{单调性}

\begin{theorem}\label{thm:calculus-monotonicity}
设\(f\)在区间\(I\)上可导。若\(f'\ge0\)，则\(f\)单调不减；若\(f'>0\)，则\(f\)严格递增。将不等号反向，得到相应的单调不增和严格递减结论。
\end{theorem}

\begin{proof}
任取\(u<v\)，由中值定理，\(f(v)-f(u)=f'(\xi)(v-u)\)。导数的符号即给出函数值之差的符号。
\end{proof}

连续端点可以一并纳入单调区间。反过来，若可导函数单调不减，则由差商定义得\(f'\ge0\)。不过，严格递增并不要求处处导数为正，例如\(x^3\)在原点导数为零。

实际判断中，若\(f'\ge0\)，且只有有限个零点，将区间在这些点处分开，再用连续性衔接，就能得到严格递增。这样既保留了函数经过水平切线的可能，也不影响不同位置之间的严格大小关系。

\begin{figure}[!ht]
\centering
\includegraphics[width=.83\linewidth]{calculus-monotonicity}
\caption{由右图导函数的符号，确定左图函数的增减；驻点两侧是否变号，决定它是否为极值点。}
\label{fig:calculus-monotonicity}
\end{figure}

\begin{example}
求\(f(x)=x^3-3x\)的单调区间。再判断\(g(x)=x^3+x\)的单调性，并与第二章不使用导数的判断比较。
\end{example}

\begin{exercise}
\mbox{}\par\nobreak

（1）求\(x e^{-x}\)与\(\ln x/x\)的单调区间。

（2）证明\(f(x)=x-\sin x\)在\(\mathbb R\)上严格递增。导数在无穷多个点为零是否妨碍结论？可在任意两点之间选取导数为正的一小段进行比较。

（3）求实数\(a\)的范围，使\(x^3-3ax\)在\(\mathbb R\)上单调不减。

（4）（2024年新课标I卷第6题改编）设
\begin{equation*}
 f(x)=\begin{cases}-x^2-2ax-a,&x<0,\\ e^x+\ln(x+1),&x\ge0.\end{cases}
\end{equation*}
若\(f\)在\(\mathbb R\)上严格递增，求\(a\)的取值范围。分段内的导数符号与连接处的取值关系，应分别考察。
\end{exercise}

\subsection{极值与最值}

\begingroup
\setlength{\parskip}{3pt}
\setlength{\abovedisplayskip}{10pt}
\setlength{\belowdisplayskip}{10pt}

\begin{proposition}\label{prop:calculus-first-derivative-test}
设\(f\)在\(x_0\)附近连续，且在\(x_0\)两侧可导。若导数在左侧为正、右侧为负，则\(x_0\)为极大值点；若左负右正，则为极小值点。若两侧同为正或同为负，则该点无极值。
\end{proposition}

这是单调性判据的直接结果：左增右减时，两侧函数值都小于中间值；其他情形类似。它给出极值的判别，而Fermat定理只给出了可导极值点必须满足的条件。

闭区间连续函数的最大值与最小值一定存在。若内部处处可导，内部最值点必为驻点，所以只须比较端点和全部驻点的函数值。驻点是候选位置，比较后才能确定最值。

以\(f(x)=x^3-3x\)在\([-3,3]\)上为例，导数在\((-3,-1)\)、\((1,3)\)上为正，在\((-1,1)\)上为负。因此\(-1\)为极大值点，\(1\)为极小值点；但
\begin{equation*}
 f(-3)=-18,\quad f(-1)=2,\quad f(1)=-2,\quad f(3)=18.
\end{equation*}
最大值为\(18\)，最小值为\(-18\)，都在端点取得。局部极值与整个区间的最值，由同一个变化过程得到，却比较不同的范围。

在开区间或无界区间上，还应研究端部行为。例如\(x+1/x\)在\(x=1\)处取得最小值\(2\)，在\((0,1)\)上递减、\((1,+\infty)\)上递增；两端均趋于正无穷，故没有最大值。

\begin{example}
求\(f(x)=x^3-3x\)在\([0,2]\)与\([0,3]\)上的最值。再求\(g(x)=\ln x-x\)在\((0,+\infty)\)上的极值和最值，比较两种范围下的研究步骤。
\end{example}

\begin{example}[2024年新课标II卷第16题]
已知\(f(x)=e^x-ax-a^3\)。
（1）当\(a=1\)时，求曲线在\((1,f(1))\)处的切线方程；
（2）若\(f\)有极小值，且极小值小于\(0\)，求\(a\)的取值范围。
\end{example}

\begin{exercise}
\mbox{}\par\nobreak

（1）求\(f(x)=x^4-2x^2\)的极值及其在\([-2,1]\)上的最值。

（2）求\(x e^{-x}\)在\([0,+\infty)\)上的最值，并判断哪些值可以取到。

（3）设\(f(x)=x^{1/3}\)。它在原点没有有限导数，是否在原点取得极值？结合\(|x|\)比较。

\begin{samepage}
（4）（2024年新课标I卷第10题）设\(f(x)=(x-1)^2(x-4)\)，判断下列结论。

（A）\(x=3\)是\(f\)的极小值点；（B）当\(0<x<1\)时，\(f(x)<f(x^2)\)；（C）当\(1<x<2\)时，\(-4<f(2x-1)<0\)；（D）当\(-1<x<0\)时，\(f(2-x)>f(x)\)。
\par
\end{samepage}
\end{exercise}
\endgroup

\begingroup
\setlength{\parskip}{3.5pt}
\setlength{\abovedisplayskip}{10pt}
\setlength{\belowdisplayskip}{10pt}
\subsection{方程的根}

研究\(f(x)=k\)，就是比较固定图像与不同高度的水平线。连续性给出中间高度的存在性，单调性控制每一段内的根数，极值则确定水平线穿过不同单调段的临界高度。

一般地，可以先将定义域分成单调区间，分别确定各段的值域，再统计\(k\)被多少段取得。公共端点只计算一次。若只需根数上界，也可以使用Rolle定理；若要确定实际根数，还应补上存在性判断。

若连续函数先严格递减、后严格递增，两端均趋于正无穷，最小值为\(m\)，则水平线\(y=k\)在\(k<m\)时不与图像相交，\(k=m\)时交于一点，\(k>m\)时左右各交于一点。单调性控制各段的交点个数，值域保证交点存在。

两端异号保证有根，两端同号却不能排除有根，如\(x^2-1\)在\([-2,2]\)上。把区间在极值点处分开后，才便于分别使用零点存在定理和单调性。

存在性证明中的取点，可以由希望得到的不等式倒推。例如，对\(a>0\)，要在\(f(x)=x+a\ln x\)中找到一个负值，可以先令\(x=e^{-t}\)、\(t>0\)，于是\(f(e^{-t})=e^{-t}-at\)。第一项小于\(1\)，所以只须使\(at\ge1\)。取\(t=1/a\)，就得到\(f(e^{-1/a})<0\)，再与\(f(1)=1\)配对即可确定一个有根区间。

这里先找容易验证的充分条件，再选择取样点。更一般地，要使\(u(x)-v(x)>0\)，可以寻找便于处理的中间数\(c\)，使\(u(x)>c>v(x)\)。这个\(c\)和取样点都由所需的符号关系决定；无需先猜到一个特别巧的数。

\begin{example}
设\(a>0\)，证明\(e^x-ax^2\)在负半轴上有且只有一个零点，并构造一个含该零点的有限区间。可以先令\(x=-t\)，利用\(e^{-t}<1\)选择使函数值为负的\(t\)。
\end{example}

\begin{example}
讨论实数\(t\)取不同值时，方程\(x^3-3x=t\)的不同实根个数，并指出根数改变对应的函数图像位置。
\end{example}

两个图像的交点问题可以写为\(f-g=0\)。有时直接研究差更自然，有时除以一个处处为正的函数，能把参数从变化的表达式中分离出来。

\begin{example}
讨论方程\(\ln x=kx\)的正实根个数。比较研究\(\ln x-kx\)与研究\(\ln x/x\)的两种组织方式。
\end{example}

\begin{exercise}
\mbox{}\par\nobreak

（1）讨论方程\(e^x=ax\)的实根个数。

（2）证明\(x^3+x=\cos x\)有且只有一个实根，并确定一个长度不超过\(1\)的有根区间。

（3）设\(f\)在开区间\(I\)上可导且\(f'\)严格递增。证明\(f\)在\(I\)内至多有两个不同零点。

（4）（2014年新课标II卷文科第21题）设\(f(x)=x^3-3x^2+ax+2\)，曲线在点\((0,2)\)处的切线与\(x\)轴交点的横坐标为\(-2\)。求\(a\)，并证明当\(k<1\)时，曲线与直线\(y=kx-2\)只有一个交点。
\end{exercise}

\subsection{不等式的参数范围}

第二章已经把恒成立与存在性条件转为最值或值域的判断。导数现在为确定这些范围提供了工具。

要求\(f(x)\ge a\)对所有\(x\in I\)成立，就是要求\(a\)不超过全部函数值。最小值\(m\)存在时，这等价于\(a\le m\)；最小值不存在时，按实际值域判断。存在性问题只要求值域中至少有一个数满足所给比较。

例如，\(x>0\)时\(\ln x/x\)的最大值为\(1/e\)，且在\(x=e\)取得。因此
\begin{equation*}
 \ln x\le ax\ \text{对所有 }x>0\text{成立}
 \quad\Longleftrightarrow\quad a\ge\frac1e.
\end{equation*}
若改为存在\(x>0\)使\(\ln x>ax\)，则条件变为\(a<1/e\)。两个问题使用同一值域，但量词决定了比较方向。

\begin{example}
求实数\(a\)的范围，使\(x^2-2ax+2\ge0\)对所有\(x\in[0,2]\)成立。比较直接研究含参数函数的最值与把参数分离后研究一个固定函数的两种方法。
\end{example}

分离参数时，乘除因子的符号和定义域决定变形是否等价。若某个位置使参数的系数为零，可在该位置先判断原条件，再在其余范围讨论。

\paragraph{先由局部条件确定必要范围} \mbox{}

若要求\(F_a(x)\ge0\)在某区间上恒成立，而\(F_a(c)=0\)，则\(c\)附近的增量也必须满足符号要求。\(c\)为内点且函数可导时，Fermat定理给出\(F_a'(c)=0\)；若只研究\(x\ge c\)，右导数则必须非负。例如，\(e^x\ge1+ax\)对所有实数成立，先在零点使用极值必要条件就得到\(a=1\)，再由\(e^x\)的单调变化或凸性验证充分性。

局部条件只是整段条件的一部分。若它尚不能确定参数，还要回到整个研究区间；后面的Taylor展开可以在低阶项抵消时进一步辨认局部增量的符号。

\paragraph{分别控制两个部分} \mbox{}

比较\(U(x)\)与\(V(x)\)，可以分别给出前者的下界、后者的上界，再比较这两个界。特别地，若存在常数\(c\)，使整个区间上\(U(x)\ge c\)、\(V(x)\le c\)，就有\(U\ge V\)。常数分离是更强的要求，一般的逐点大小关系并不保证\(\min U\ge\max V\)。

例如，在\(x>0\)时，由\(e^x\ge1+x\)、\(\ln x\le x-1\)，相减就得到\(e^x-\ln x\ge2\)。选择中间表达式\(x+1\)、\(x-1\)，使共同的\(x\)消去，比直接寻找\(e^x-\ln x\)的最小值更方便。所选估计的方向与消去后留下的量，需要一同考虑。

当条件中出现两个自变量时，还应先确定它们怎样取值。设\(F,G\)分别在闭区间\(I,J\)上连续，记各自的最大值为\(M_F,M_G\)，\(G\)的最小值为\(m_G\)。则
\begin{equation*}
 \begin{aligned}
 &\text{对每个 }x\in I\text{，存在 }y\in J\text{ 使 }F(x)\le G(y)
 &&\Longleftrightarrow M_F\le M_G,\\
 &\text{对每个 }x\in I\text{ 和每个 }y\in J\text{ 都有 }F(x)\le G(y)
 &&\Longleftrightarrow M_F\le m_G.
 \end{aligned}
\end{equation*}
前者允许选择\(G\)的最大值与全部\(F(x)\)比较，后者则连\(G\)的最小值也必须满足要求。导数负责求出这些最值，量词决定应当比较哪两个最值。

\begin{exercise}
\mbox{}\par\nobreak

（1）求\(a\)的范围，使\(\ln x\le a(x-1)\)对所有\(x>0\)成立。

（2）分别求\(a\)的范围，使\(e^x-ax\ge0\)对所有实数\(x\)成立，以及使存在实数\(x\)满足\(e^x-ax<0\)。

（3）求\(t\)的范围，使对每个\(x\in[1,e]\)，都存在\(y\in[0,1]\)满足\(\ln x-t\le y-y^2\)。
\end{exercise}

\subsection{含参数函数的讨论}

含参数函数的分类，应依据单调结构、候选极值点或根数发生改变的位置。常见的临界情形包括导数零点合并、驻点进入研究区间、极值与指定高度相等。

例如，对\(f_a(x)=x^3-3ax\)，导数为\(3(x^2-a)\)。当\(a\le0\)时，函数严格递增；当\(a>0\)时，出现两个驻点\(\pm\sqrt a\)，对应的极大值、极小值分别为\(2a^{3/2}\)、\(-2a^{3/2}\)。因此讨论方程\(f_a(x)=t\)，可以先按\(a\)判断图像结构，再比较\(t\)与这两个临界高度。

若只研究\([0,1]\)，还需要比较驻点与端点的位置。参数\(a=1\)的意义来自\(\sqrt a=1\)，而不是因为某个代数表达式恰好可以按这个数分类。

% 题号按本地试卷PDF核对；转录汇编的条目编号与试卷有一题偏移。
\begin{example}[2023年新课标II卷第11题]
若函数\(f(x)=a\ln x+b/x+c/x^2\)（\(a\ne0\)）既有极大值也有极小值，则下列哪些结论正确？
\begin{equation*}
 \text{A. }bc>0\qquad \text{B. }ab>0\qquad
 \text{C. }b^2+8ac>0\qquad \text{D. }ac<0.
\end{equation*}
\end{example}

这类问题的导数往往可以写成一个正因子乘较简单的表达式。先确认符号由哪一部分决定，再研究其零点位置和变号情况，可以把参数条件与函数结构对应起来。

导数方程的实根还要放回原定义域。例如在\(x>0\)上，负根不会产生驻点；两个正根合并时，导数也未必变号。因此判别式、根的位置和两侧符号分别回答不同的问题，不能仅凭导数方程有两个实根就断定原函数有两个极值。

\begin{exercise}\label{exs:calculus-parameter-gaokao}
\mbox{}\par\nobreak

（1）讨论\(f_a(x)=x^3-3ax\)在\([0,1]\)上的最值，并指出需要分类的参数位置。

（2）设\(f_a(x)=\ln x-ax\)。讨论其单调性，并判断不同\(a\)下\(f_a(x)=0\)的根数。

（3）（2024年新课标I卷第18题）已知\(f(x)=\ln\frac{x}{2-x}+ax+b(x-1)^3\)。

\exercisesubitem{1}若\(b=0\)且\(f'(x)\ge0\)，求\(a\)的最小值；

\exercisesubitem{2}证明曲线\(y=f(x)\)是中心对称图形；

\exercisesubitem{3}若\(f(x)>-2\)当且仅当\(1<x<2\)，求\(b\)的取值范围。
\end{exercise}

\subsection{切线问题}

给定切点时，直接使用式\eqref{eq:calculus-tangent}。若给定的是切线经过的点，就应把切点作为未知量。

设切点横坐标为\(t\)，切线经过\((u,v)\)，则
\begin{equation}\label{eq:calculus-tangent-through-point}
 v=f(t)+f'(t)(u-t).
\end{equation}
原问题由此变为关于\(t\)的方程。\((u,v)\)未必是切点，不能直接用\(f'(u)\)代替切线斜率。

\begin{example}
求过点\((0,-1)\)且与\(y=x^2\)相切的所有直线。再求过\((0,b)\)且与\(y=\ln x\)相切的直线，并讨论其条数。
\end{example}

设两条曲线为\(y=f(x)\)、\(y=g(x)\)，切点横坐标分别为\(s,t\)。令两条切线的斜率与截距分别相同，得到公共切线的条件：
\begin{equation*}
 f'(s)=g'(t),\qquad f(s)-sf'(s)=g(t)-tg'(t).
\end{equation*}
两切点可以不同。若要求在公共点相切，还须\(s=t\)且\(f(s)=g(t)\)。几何约束可再联用第四章的曲线方程与向量方法。

\begin{exercise}
\mbox{}\par\nobreak

（1）求\(y=e^x\)上平行于\(y=2x\)的切线。

（2）求\(y=x^2\)与\(y=-(x-2)^2\)的公共切线。

（3）求实数\(a\)，使\(y=\ln x\)与\(y=ax^2\)有公共点，且在该公共点处切线相同。
\end{exercise}

\subsection{最优化问题}

导数不能替代建模。先确定变化的量、约束及允许范围，再把目标量表示成单变量函数，才能使用最值方法。最后还应把候选点还原到原问题，核对它是否满足约束。

例如，用长为\(L\)的材料围成矩形，一边为\(x\)，另一边为\(L/2-x\)，面积为\(A(x)=x(L/2-x)\)，其中\(0<x<L/2\)。由\(A'(x)=L/2-2x\)，面积先增后减，在\(x=L/4\)处取得最大值。这个结论也可以用配方或均值不等式得到；导数把同一种研究推进到更一般的目标函数。

\begin{example}
在半径为\(R\)的圆内作内接矩形，求面积的最大值。若改为在半轴长为\(a,b\)的椭圆内作边平行于坐标轴的内接矩形，怎样结合伸缩变换与导数研究？
\end{example}

\begin{example}
从边长为\(L\)的正方形纸板四角各剪去一个边长为\(x\)的小正方形，再折成无盖长方体盒。建立体积关于\(x\)的函数，求体积最大时的\(x\)。
\end{example}

约束使目标量同时受到几个相反方向的影响：盒子变高时，底面也在缩小。把它们写进同一个函数，才能比较总体变化。结果还应符合原量纲：所有长度放大为原来的\(k\)倍时，面积与体积分别放大为\(k^2\)、\(k^3\)倍。

\begin{exercise}
\mbox{}\par\nobreak

（1）制作一个体积为\(V>0\)的无盖圆柱容器，侧面与底面的单位面积成本相同。求用料最少时的底面半径与高的关系。

（2）在曲线\(y=1/x\)的第一象限部分上取点\(P\)，使\(P\)到原点的距离最小。求\(P\)，并解释为什么可以改为研究距离的平方。

（3）设\(a>0\)，点\(A=(0,a)\)、\(B=(1,0)\)。在\(x\)轴上取\(P=(t,0)\)，\(0\le t\le1\)，求\(|AP|+\tfrac12|PB|\)的最小值。驻点不在研究区间内时，应怎样完成判断？
\end{exercise}
\endgroup
